% GENERATED FILE. Edit canonical lessons, then run npm run generate:books.
% canonical-source-hash: fnv1a64:9f7f7df2027cc59d
% canonical-lessons: GU-C432-sonu, GU-C432-cadi, GU-C432-jhaverat, GU-C432-ghasvu, GU-C432-camkavvu, GU-R432-review-425-432

\chapter{Gold and Silver}
\label{ch:gu-gold-and-silver}
\hlchaptermodality{432}

\begin{chapteropening}
\textbf{By the end of this chapter:} \emph{I can talk about gold, silver and jewellery.}

Complete the last lesson of chapter 432: I can talk about gold, silver and jewellery.
\end{chapteropening}

\section[\texorpdfstring{\gu{સોનું} (sonũ)}{સોનું (sonũ)}]{\gu{સોનું} (sonũ) — gold}

\label{lesson:GU-C432-sonu}



\begin{quote}
Before the new one: say the Gujarati for wholesale, then the Gujarati for to hoard.
\end{quote}

\subsection*{You'll want to know: \gu{સોનું}}
\textbf{\gu{સોનું}} — \emph{sonũ} — \textquotedblleft{}gold\textquotedblright{}.

\begin{quote}
\gu{દાદી} \gu{પાસે} \gu{સોનાની} \gu{વીંટી} \gu{છે}.
\end{quote}

\begin{quote}
\emph{dādī pāse sonānī vī̃ṭī che.}
\end{quote}

\begin{quote}
Grandmother has a gold ring.
\end{quote}

A neuter noun.

\begin{cousinweb}
\textbf{\gu{સોનું}} (\emph{sonũ}) is a tadbhava of Sanskrit \emph{suvarṇa}, \textquotedblleft{}of a good colour\textquotedblright{}, and so \textquotedblleft{}gold\textquotedblright{}; CDIAL lists the Gujarati word (CDIAL 13519).
\end{cousinweb}

\begin{grammarlens}[title={What you've built}]
One more word for getting around.
\end{grammarlens}

\subsection*{Your turn}
\begin{itemize}
\raggedright
  \item \emph{Say it:} \emph{sonũ}
  \item \emph{Say it:} \emph{sonũ}, once more
  \item \emph{Say it:} \emph{dādī pāse sonānī vī̃ṭī che.} — Grandmother has a gold ring.
  \item \emph{Recall:} say the Gujarati for wholesale, then the Gujarati for to hoard, then say \emph{sonũ} again
\end{itemize}

\begin{tcolorbox}[breakable,colback=teal!4,colframe=teal!35!black,title={Before you move on}]
What does \gu{સોનું} mean? (\textquotedblleft{}Gold\textquotedblright{}.) Say it once more.
\end{tcolorbox}
\section[\texorpdfstring{\gu{ચાંદી} (cā̃dī)}{ચાંદી (cā̃dī)}]{\gu{ચાંદી} (cā̃dī) — silver}

\label{lesson:GU-C432-cadi}



\begin{quote}
Before the new one: say the Gujarati for to hoard, then the Gujarati for gold.
\end{quote}

\subsection*{You'll want to know: \gu{ચાંદી}}
\textbf{\gu{ચાંદી}} — \emph{cā̃dī} — \textquotedblleft{}silver\textquotedblright{}.

\begin{quote}
\gu{આ} \gu{વાટકી} \gu{ચાંદીની} \gu{છે}.
\end{quote}

\begin{quote}
\emph{ā vāṭkī cā̃dīnī che.}
\end{quote}

\begin{quote}
This little bowl is made of silver.
\end{quote}

A feminine noun.

\begin{cousinweb}
\textbf{\gu{ચાંદી}} (\emph{cā̃dī}) is a tadbhava of Sanskrit \emph{candrikā}, \textquotedblleft{}moonlight\textquotedblright{}; CDIAL lists the Gujarati word, \textquotedblleft{}silver\textquotedblright{} (CDIAL 4669): the metal with the moon's colour.
\end{cousinweb}

\begin{grammarlens}[title={What you've built}]
One more word for getting around.
\end{grammarlens}

\subsection*{Your turn}
\begin{itemize}
\raggedright
  \item \emph{Say it:} \emph{cā̃dī}
  \item \emph{Say it:} \emph{cā̃dī}, once more
  \item \emph{Say it:} \emph{ā vāṭkī cā̃dīnī che.} — This little bowl is made of silver.
  \item \emph{Recall:} say the Gujarati for to hoard, then the Gujarati for gold, then say \emph{cā̃dī} again
\end{itemize}

\begin{tcolorbox}[breakable,colback=teal!4,colframe=teal!35!black,title={Before you move on}]
What does \gu{ચાંદી} mean? (\textquotedblleft{}Silver\textquotedblright{}.) Say it once more.
\end{tcolorbox}
\section[\texorpdfstring{\gu{ઝવેરાત} (jhaverāt)}{ઝવેરાત (jhaverāt)}]{\gu{ઝવેરાત} (jhaverāt) — jewellery; gems}

\label{lesson:GU-C432-jhaverat}



\begin{quote}
Before the new one: say the Gujarati for gold, then the Gujarati for silver.
\end{quote}

\subsection*{You'll want to know: \gu{ઝવેરાત}}
\textbf{\gu{ઝવેરાત}} — \emph{jhaverāt} — \textquotedblleft{}jewellery; gems\textquotedblright{}.

\begin{quote}
\gu{લગ્ન} \gu{માટે} \gu{તેમણે} \gu{ઝવેરાત} \gu{ખરીદ્યું}.
\end{quote}

\begin{quote}
\emph{lagna māṭe temṇe jhaverāt kharīdyũ.}
\end{quote}

\begin{quote}
They bought jewellery for the wedding.
\end{quote}

A neuter noun.

\begin{cousinweb}
\textbf{\gu{ઝવેરાત}} (\emph{jhaverāt}) is usually traced to Arabic \emph{jawāhirāt}, \textquotedblleft{}jewels\textquotedblright{}, which came through Persian.
\end{cousinweb}

\begin{grammarlens}[title={What you've built}]
One more word for getting around.
\end{grammarlens}

\subsection*{Your turn}
\begin{itemize}
\raggedright
  \item \emph{Say it:} \emph{jhaverāt}
  \item \emph{Say it:} \emph{jhaverāt}, once more
  \item \emph{Say it:} \emph{lagna māṭe temṇe jhaverāt kharīdyũ.} — They bought jewellery for the wedding.
  \item \emph{Recall:} say the Gujarati for gold, then the Gujarati for silver, then say \emph{jhaverāt} again
\end{itemize}

\begin{tcolorbox}[breakable,colback=teal!4,colframe=teal!35!black,title={Before you move on}]
What does \gu{ઝવેરાત} mean? (\textquotedblleft{}Jewellery; gems\textquotedblright{}.) Say it once more.
\end{tcolorbox}
\section[\texorpdfstring{\gu{ઘસવું} (ghasvũ)}{ઘસવું (ghasvũ)}]{\gu{ઘસવું} (ghasvũ) — to rub; to polish}

\label{lesson:GU-C432-ghasvu}



\begin{quote}
Before the new one: say the Gujarati for silver, then the Gujarati for jewellery.
\end{quote}

\subsection*{You'll want to know: \gu{ઘસવું}}
\textbf{\gu{ઘસવું}} — \emph{ghasvũ} — \textquotedblleft{}to rub; to polish\textquotedblright{}.

\begin{quote}
\gu{તેણે} \gu{ચાંદીની} \gu{વાટકી} \gu{કપડાથી} \gu{ઘસી}.
\end{quote}

\begin{quote}
\emph{teṇe cā̃dīnī vāṭkī kapḍāthī ghasī.}
\end{quote}

\begin{quote}
She rubbed the silver bowl with a cloth.
\end{quote}

Past, with \gu{તેણે} (he or she …): \gu{તેણે} \gu{વાસણ} \gu{ઘસ્યું}, \emph{teṇe vāsaṇ ghasyũ}.

\begin{cousinweb}
\textbf{\gu{ઘસવું}} (\emph{ghasvũ}) is a tadbhava of Sanskrit \emph{gharṣati}, \textquotedblleft{}rubs\textquotedblright{}; CDIAL lists the Gujarati verb, \textquotedblleft{}to rub, to polish\textquotedblright{} (CDIAL 4450).
\end{cousinweb}

\begin{grammarlens}[title={What you've built}]
One more word for getting around.
\end{grammarlens}

\subsection*{Your turn}
\begin{itemize}
\raggedright
  \item \emph{Say it:} \emph{ghasvũ}
  \item \emph{Say it:} \emph{ghasvũ}, once more
  \item \emph{Say it:} \emph{teṇe cā̃dīnī vāṭkī kapḍāthī ghasī.} — She rubbed the silver bowl with a cloth.
  \item \emph{Recall:} say the Gujarati for silver, then the Gujarati for jewellery, then say \emph{ghasvũ} again
\end{itemize}

\begin{tcolorbox}[breakable,colback=teal!4,colframe=teal!35!black,title={Before you move on}]
What does \gu{ઘસવું} mean? (\textquotedblleft{}To rub; to polish\textquotedblright{}.) Say it once more.
\end{tcolorbox}
\section[\texorpdfstring{\gu{ચમકાવવું} (camkāvvũ)}{ચમકાવવું (camkāvvũ)}]{\gu{ચમકાવવું} (camkāvvũ) — to make shine; to polish until it shines}

\label{lesson:GU-C432-camkavvu}



\begin{quote}
Before the new one: say the Gujarati for jewellery, then the Gujarati for to rub.
\end{quote}

\subsection*{You'll want to know: \gu{ચમકાવવું}}
\textbf{\gu{ચમકાવવું}} — \emph{camkāvvũ} — \textquotedblleft{}to make shine; to polish until it shines\textquotedblright{}.

\begin{quote}
\gu{તેણે} \gu{જૂનું} \gu{ઘરેણું} \gu{ઘસીને} \gu{ચમકાવ્યું}.
\end{quote}

\begin{quote}
\emph{teṇe jūnũ ghareṇũ ghasīne camkāvyũ.}
\end{quote}

\begin{quote}
He rubbed the old ornament until it shone.
\end{quote}

Past, with \gu{તેણે} (he or she …): \gu{તેણે} \gu{વીંટી} \gu{ચમકાવી}, \emph{teṇe vī̃ṭī camkāvī}.

\begin{cousinweb}
\textbf{\gu{ચમકાવવું}} (\emph{camkāvvũ}) is the causative partner of \gu{ચમકવું}, \textquotedblleft{}to shine\textquotedblright{}, which goes back to \emph{*camakka}, a form CDIAL reconstructs (CDIAL 4676): to make a thing shine.
\end{cousinweb}

\begin{grammarlens}[title={What you've built}]
One more word for getting around.
\end{grammarlens}

\subsection*{Your turn}
\begin{itemize}
\raggedright
  \item \emph{Say it:} \emph{camkāvvũ}
  \item \emph{Say it:} \emph{camkāvvũ}, once more
  \item \emph{Say it:} \emph{teṇe jūnũ ghareṇũ ghasīne camkāvyũ.} — He rubbed the old ornament until it shone.
  \item \emph{Recall:} say the Gujarati for jewellery, then the Gujarati for to rub, then say \emph{camkāvvũ} again
\end{itemize}

\begin{tcolorbox}[breakable,colback=teal!4,colframe=teal!35!black,title={Before you move on}]
What does \gu{ચમકાવવું} mean? (\textquotedblleft{}To make shine; to polish until it shines\textquotedblright{}.) Say it once more.
\end{tcolorbox}
\section[(dialogue)]{Looking back: from Loans and Instalments to Gold and Silver}

\label{lesson:GU-R432-review-425-432}



\begin{quote}
Say the Gujarati for a bank loan and an instalment.
\end{quote}

\subsection*{Your turn}
Say these aloud:

\begin{itemize}
\raggedright
  \item a bank loan, an instalment, a broker, a pledge, paid in full
  \item a lakh, a crore, subtraction, to multiply, to increase
  \item a pension, an investment, a deposit, a fixed term, to allot
  \item bankruptcy, money owed to you, borrowed for a while, a moneylender, to swindle
  \item a safe, funds, a currency, to be robbed, to be affordable
  \item a pair of scales, a maund, a small packet, a regular buyer, to weigh out
  \item a warehouse, a supply, a shortfall, wholesale, to hoard
  \item gold, silver, jewellery, to rub, to make shine
\end{itemize}

\begin{tcolorbox}[breakable,colback=teal!4,colframe=teal!35!black,title={Before you move on}]
A bank loan? (\textbf{\gu{લોન}}, \emph{lon}.) An instalment? (\textbf{\gu{હપતો}}, \emph{hapto}.) A broker? (\textbf{\gu{દલાલ}}, \emph{dalāl}.) A pledge? (\textbf{\gu{ગીરો}}, \emph{gīro}.) Paid in full? (\textbf{\gu{ચૂકતે}}, \emph{cūkte}.) A lakh? (\textbf{\gu{લાખ}}, \emph{lākh}.) A crore? (\textbf{\gu{કરોડ}}, \emph{karoḍ}.) Subtraction? (\textbf{\gu{બાદબાકી}}, \emph{bādbākī}.) To multiply? (\textbf{\gu{ગુણવું}}, \emph{guṇvũ}.) To increase? (\textbf{\gu{વધારવું}}, \emph{vadhārvũ}.) A pension? (\textbf{\gu{પેન્શન}}, \emph{penśan}.) An investment? (\textbf{\gu{રોકાણ}}, \emph{rokāṇ}.) A deposit? (\textbf{\gu{થાપણ}}, \emph{thāpaṇ}.) A fixed term? (\textbf{\gu{મુદત}}, \emph{mudat}.) To allot? (\textbf{\gu{ફાળવવું}}, \emph{phāḷavvũ}.) Bankruptcy? (\textbf{\gu{દેવાળું}}, \emph{devāḷũ}.) Money owed to you? (\textbf{\gu{લેણું}}, \emph{leṇũ}.) Borrowed for a while? (\textbf{\gu{ઉછીનું}}, \emph{uchīnũ}.) A moneylender? (\textbf{\gu{શાહુકાર}}, \emph{śāhukār}.) To swindle? (\textbf{\gu{ઠગવું}}, \emph{ṭhagvũ}.) A safe? (\textbf{\gu{તિજોરી}}, \emph{tijorī}.) Funds? (\textbf{\gu{નાણાં}}, \emph{nāṇā̃}.) A currency? (\textbf{\gu{ચલણ}}, \emph{calaṇ}.) To be robbed? (\textbf{\gu{લૂંટાવું}}, \emph{lū̃ṭāvũ}.) To be affordable? (\textbf{\gu{પરવડવું}}, \emph{parvaḍvũ}.) A pair of scales? (\textbf{\gu{ત્રાજવું}}, \emph{trājvũ}.) A maund? (\textbf{\gu{મણ}}, \emph{maṇ}.) A small packet? (\textbf{\gu{પડીકું}}, \emph{paḍīkũ}.) A regular buyer? (\textbf{\gu{ઘરાક}}, \emph{gharāk}.) To weigh out? (\textbf{\gu{જોખવું}}, \emph{jokhvũ}.) A warehouse? (\textbf{\gu{ગોદામ}}, \emph{godām}.) A supply? (\textbf{\gu{પુરવઠો}}, \emph{puravṭho}.) A shortfall? (\textbf{\gu{ઊણપ}}, \emph{ūṇap}.) Wholesale? (\textbf{\gu{જથ્થાબંધ}}, \emph{jaththābandh}.) To hoard? (\textbf{\gu{સંઘરવું}}, \emph{sangharvũ}.) Gold? (\textbf{\gu{સોનું}}, \emph{sonũ}.) Silver? (\textbf{\gu{ચાંદી}}, \emph{cā̃dī}.) Jewellery? (\textbf{\gu{ઝવેરાત}}, \emph{jhaverāt}.) To rub? (\textbf{\gu{ઘસવું}}, \emph{ghasvũ}.) To make shine? (\textbf{\gu{ચમકાવવું}}, \emph{camkāvvũ}.)
\end{tcolorbox}
